\documentclass{article}

% Language setting
% Replace `english' with e.g. `spanish' to change the document language
\usepackage[english]{babel}

% Set page size and margins
% Replace `letterpaper' with `a4paper' for UK/EU standard size
\usepackage[a4paper,top=2cm,bottom=2cm,left=2cm,right=2cm,marginparwidth=0cm]{geometry}

% Useful packages
\usepackage{amsmath}
\usepackage{graphicx}
\usepackage[colorlinks=true, allcolors=blue]{hyperref}
\usepackage{float}
\title{General Assessment}
\author{Ahmed Homoudi} %and Klemens Barfus
\usepackage{setspace}
\onehalfspacing 
\usepackage{caption}
\usepackage{subcaption}
\begin{document}
\maketitle
\noindent 


\noindent
\textbf{Remarks:} 
\begin{itemize}
\item Use any programming languages you would like to. Keep in mind that it should be a language capable of handling scientific computations. We, generally, are familiar with R and Python.
\item The Internet is filled with Q\&A regarding programming, and there is definitely someone who tried to program what you want to do. So, Google is your friend.
\item Please submit your code and results in a readable manner. 
\item Comment your code, and leave notes, therefore, you or other people can read it easily later. 
\end{itemize}

\section{Task 1: String manipulation}
Write a code that can do these tasks: 

\begin{itemize}
\item Print "Hello World!". 
\item Check out your working directory (WD) and Print it.
\item Make a new directory from within your code. 
\item Change your WD to the newly created directory. 
\item Create a TXT file in the WD, and Write in it: \\
               "Hello world!"\\
               "I'm your name"
\item Read that text file and replace "World" with "Dresden" or any other word, and finally write it to the TXT file. 
    
\item Advanced task: If you want to challenge yourself, this is the chance: 
    \begin{itemize}
    \item Create a vector of string that has a length of 50 elements. Each element is, for example, "Hello! I'm file number n". n is the order of the string in the vector. 
    \item Loop over the vector elements, and write each element to a TXT file called "file\_n.txt"
    \item List all these TXT files and read their content to one variable, and write this variable to a TXT file called "all\_n.txt"
    \item Delete the other 50 files. 
    \end{itemize}
\end{itemize}

\section{Task 2: One Dimensional Data}
Go to \href{https://climateknowledgeportal.worldbank.org/download-data}{Climate Change Knowledge Portal}. Download a time series data of precipitation and mean temperature for your country or area of interest. After downloading the data as shown in Figure \ref{fig001}, continue in the same code you started in Task 1 and write a script that: 

\begin{figure}[H]
         
    \centering
    \includegraphics[width=\textwidth]{Screenshot from 2022-08-05 10-27-08.png}
    \caption{Climate Change Knowledge Portal. CRU (Observed)}
    \label{fig001}
    

\end{figure}

\begin{itemize}
\item Read the two CSV file you downloaded correctly. 
\item Plot a time series of the two variables.
\item Produce a long-term climatology from the data. Look at the first left plot \href{https://en.climate-data.org/europe/bulgaria/german/german-198333/}{here}. %If you would like a challenge, include both datasets in the same graph. 
    
\item Advanced task: If you want to challenge yourself, this is another chance: 
    \begin{itemize}
    \item calculate the yearly accumulated rainfall and plot it as bar-plot where x-axis is years. 
    \item add a trend line with slope value on the chart. 
    \end{itemize}
\end{itemize}

\section{Task 3: Two Dimensional Data}

Go to \href{https://cloudstore.zih.tu-dresden.de/index.php/s/Bks4sAZbPD46SZ7}{CloudStore, TU Dresden} and download the file \textit{task3\_file.nc}. The file contains 2D data (excluding time dimension) from ERA5 reanalysis including 2 m temperature, total precipitation, u and v wind compounds. After downloading the file, continue in the same code you started in tasks above and write a script that: 

\begin{itemize}
\item Read the netCDF file.
\item Extract different variables from the file. 
\item Determine the spatial resolution of data in longitude and latitude directions. 
\item search for the grid cell that represents your region of interest, and extract the variables as monthly time series. 
\item Plot 2 m temperature, total precipitation in one time series plot (Hint: two different y-axis).
\item Plot the u and v wind components as well as wind speed. 
    
\item Advance task: If you want to challenge yourself, this is another chance: 
    \begin{itemize}
    \item Aggregate the variables on the time axis and produce a long term climatology. 
    \item Plot a map focused on your region showing long-term climatology of 2 m temperature, total precipitation, and wind. (Hint: the data has dimension of (time=264, lat=180, lon=360), a worldwide climatology should have (time=1, lat=180, lon=360)).
    \item Try to add map legends, shape files indicating the borders and every features that is part of map standards.
    \item You could plot wind speed and directions as arrows.  
    \end{itemize}
\end{itemize}



\section{Task 4: Three Dimensional Data (Advanced)}
Go to \href{https://cloudstore.zih.tu-dresden.de/index.php/s/Bks4sAZbPD46SZ7}{CloudStore, TU Dresden} and download the file \textit{task4\_file.nc}. The file contains 3D data (excluding time dimension) from ERA5 reanalysis including temperature, relative humidity, geopotential (See more about it \href{https://apps.ecmwf.int/codes/grib/param-db/?id=129}{here}), and u and v wind compounds. After downloading the file, continue in the same code you started in tasks above and write a script that: 

\begin{itemize}
\item Read the netCDF file.
\item Extract different variables from the file. 
\item Determine the spatial resolution of data in longitude and latitude directions. 
\item Search for the grid cell that represents your region of interest, and extract the variables as monthly time series. 
    (Hint: the file contains atmospheric data on variable levels for one year. So the monthly time series should a table for the location you picked containing 3 columns: month, level, value)
\item Plot atmospheric profiles for different variable, where x-axis is the climatic variable, y-axis is pressure level or/and altitude. Start with plotting one month. Check \href{https://en.wikipedia.org/wiki/U.S._Standard_Atmosphere}{the first figure on the right there}. 
\item Plot similar profile the u and v wind components as well as wind speed. 
    
\item Advanced task: If you want to challenge yourself, this is  another chance: 
    \begin{itemize}
    \item In the plots above, add a profile line for each month for each variable. 
    \end{itemize}
\end{itemize}


\vspace{5mm}
\begin{center}
    \textbf{End}
\end{center}

\newpage
\section{Internal: New supervision approach}
After the student show an interest in the topic, they will receive the 'General assessment test'. They will have 1 to 2 weeks to solve the task above. Afterwards, they should submit their code, their actual grades until now, and a detailed CV. 

\begin{itemize}
\item No matrix chat, email only. It should be one thread; including all supervisors.  
\item A short, consistent email about the progress and issues weekly.
\item Meetings will be held triweekly. Preferably, on Friday. In case of urgent issues, a meeting can be organised in the following week.
\item Meetings require the presence of all supervisors. 
\item Emails during the weekend are unlikely to be answered.
\end{itemize}

%\section{General Information}

%\section{Data Knoweldge} 
%\begin{enumerate}
%\item If you have being supplied with a CSV file including daily temperature data for your home town. Can you elaborate what you expect in this file? 
%\end{enumerate}

%\section{Climatological Knowledge} 
%\textbf{Attention:} In this section, the applicant is asked to present their knowledge about climate and weather process. It will not affect the general assessment score. 
%Please try to answer the questions here without external help. There is no issue in googling the terms or using a dictionary. 
%timeseries 
%maps 
%2d time series 
%3D time series 
%Radisende 
%clouds 
%general circulation 
%temperature, humidity, wind 
%lat and lon 
%describe the climate at your hometown considering temperature and 
%\begin{Form}
%\TextField{Last Name:} \\
%\TextField{First Name:} \\
%\TextField{Date:} \\
%\TextField{Course of Study:} \\
%\end{Form}
\end{document}
